\documentclass[10pt,a4paper]{article}
\usepackage[T1]{fontenc}
\usepackage[utf8]{inputenc}
\usepackage[french]{babel}
\usepackage{graphicx,fancyhdr,}
\usepackage{mathtools}
%\usepackage[left=1cm,right=1cm,top=1.5cm,bottom=1.5cm]{geometry}
\usepackage{amsfonts,amsmath,amssymb}
\usepackage{amsthm}
\usepackage{pstricks-add, pspicture} % pour les axces
\usepackage{thmtools}
\usepackage{xcolor}
\usepackage{nameref}
\usepackage{enumitem}
\usepackage{hyperref}
\usepackage{hyperref, fancybox, microtype,pifont,array, caption, bbding, multicol, longtable, indentfirst, tabularx, ,setspace, fourier, courier, color, multirow, lscape,lipsum,textcomp}
\usepackage{tikz}
\usepackage{systeme}
\usepackage{pgfplots}
\pgfplotsset{compat=1.15}
\usepackage{mathrsfs}
\usetikzlibrary{arrows}
\pagestyle{empty}


%<<<<<<<WARNING>>>>>>>
% PGF/Tikz doesn't support hatch filling very well
% Use PStricks for a perfect hatching export

\usetikzlibrary[patterns]

\usepackage[normalem]{ulem} % pour hachurer
\usepackage{colortbl} %Pour hachurer une case d'un tableau 
\usepackage[tikz]{bclogo}
\usepackage[most]{tcolorbox} 
\usepackage{longtable}
%\usepackage{background} %filigrane
\usepackage{draftwatermark} %filigrane
\usepackage{fancybox,fancyhdr}%encadrement et les titres des en-tete
%\usepackage{multicol, multirow}%fusion de colone et ligne
\usepackage[np]{numprint}%pour écrire correctement les nombres à virgules sans espaces entre la virgule et le nombre tout en utilisant le commande \np{}
%\usepackage{bookman, fourier}

%\setlength{\columnseprule}{0.5pt}
\setlength{\columnsep}{3pt}
\usepackage{graphicx}
\usepackage[left=1.00cm, right=1.00cm, top=1.2cm, bottom=1.2cm]{geometry}
\newcounter{exo}
\setcounter{exo}{0}
\newcommand{\exo}[1]{
\stepcounter{exo} { \textbf{\underline{Exercice de maison \No}\arabic{exo}} : \textbf{#1} \newline }}
%donnez une fraction aleatoire à simplifier debut:
\newcounter{Num} \newcounter{Den}
\newcounter{Coef}
\newcommand{\FractAleat}{%
	\reinitrand[first=1,last=10,counter=Num]\rand
	\reinitrand[first=2,last=10,counter=Den]\rand
	\reinitrand[first=2,last=12,counter=Coef]\rand
	\setcounter{Num}{\value{Num}*\value{Coef}}
	\setcounter{Den}{\value{Den}*\value{Coef}}
	$\dfrac{\theNum}{\theDen}$}
%fin
%les activités et les consignes debut
\newcounter{q}
\setcounter{q}{0}
\newcounter{qq}
\newcommand{\act}[1]
{ 
		 \setcounter{qq}{0}
	\addtocounter{q}{1} \doublebox {\textbf{Activité} \theq.}\space \textbf{{#1}}\\
 }
\newcommand{\conAct}[1]{
	\addtocounter{qq}{1} \ovalbox{\textbf{Consignes	\theq.\space\theqq.}}  \textbf{{#1}} \\
}
\newcommand{\resul}{ \begin{center}
		\rule{2cm}{1mm} 
		\shadowbox{Résolutions}
  \rule{2cm}{1mm}
\end{center}}
%definition et autre 
\newcommand{\defs}[1]{\begin{bclogo}[couleur=green!20,arrondi=5pt]
		{\textbf{Définition}}
		#1
	\end{bclogo} 	 
}
%nb ou vocabulaire 
\newcommand{\nbVo}[2]{\begin{bclogo}[couleur = lightgray!15, couleur =red!15, arrondi = 0.3, ombre = true, epOmbre =0.25, couleurOmbre= lightgray!90, logo=\bcrecyclage, epBarre=3  ]{\textbf{#1}}
		#2
\end{bclogo}}
%Propriété
\newcommand{\propr}[1]{  \begin{bclogo}[couleur = lightgray!15, arrondi = 0.3, ombre = true, epOmbre =0.25,  nobreak=false, couleurOmbre= lightgray!90, logo=\bcrecyclage, epBarre=3 ]{Propriété} #1
\end{bclogo}}
%suitepro
\newcommand{\prosuit}[1]{  \begin{bclogo}[couleur = lightgray!15, arrondi = 0.3, ombre = true, epOmbre =0.25, couleurOmbre= lightgray!90, logo=\bcrecyclage,  epBarre=3 ]{   } #1
\end{bclogo}}
%attention 
\newcommand{\atten}[1]{  \begin{bclogo}[couleur = red!25, arrondi = 0.3, ombre = true, epOmbre =0.25, couleurOmbre= lightgray!90, logo=\bcattention, epBarre=3 ]{\textcolor{red!100}{ATTENTION}} #1
\end{bclogo}}
%les retenons 
\newcommand{\ret}[1]{\begin{tcolorbox}[colback=violet!10,colframe=blue!40,title=Retenons]
		#1 
\end{tcolorbox}}
	\newcommand{\dure}{ \underline{Stratégies et Durées:} TI : \dots min TG: \dots min   TC : \dots min \\
	\ding{36}\dotfill{} }
%titre des sequence 
	\newcommand{\dures}{ \underline{Stratégies et Durées:} TI : \dots min TG: \dots min   TC : \dots min \\
	\ding{36}\dotfill{} }
\newcounter{sequen}
\setcounter{sequen}{0}
\newcommand{\sequence}[1]{
	\stepcounter{sequen}{
		\begin{bclogo}[couleur=cyan!50,arrondi=0.3, ombre = true, epOmbre =0.25,logo=\bcsoleil,epBarre=3.5]{ \LARGE Séquence \no\thesequen : }
			\Large \textbf{#1}
\end{bclogo} } }
	\newcounter{nu}
\setcounter{nu}{0}
\newcommand{\seq}{\addtocounter{nu}{0} 
	\textbf{\fcolorbox{red}{yellow}{{\LARGE Séquence} \no\thenu} \hspace{0.5cm}}}
%titre de la SA
\newcounter{situa}
\setcounter{situa}{0}
\newcommand{\situ}[1]{
	\stepcounter{situa}{
 \begin{bclogo}[couleur=blue!30,arrondi=0.3, ombre = true, epOmbre =0.25,logo=\bcbook,epBarre=3.5]{ \LARGE Situation d'Apprentissage \no\thesitua : }
 	\Large \textbf{#1}
\end{bclogo} } }
\setlength{\columnsep}{30pt}
%inclu
\newcommand{\nincl}{$\subset  $ \hspace*{-0.4cm}/ }
\newcommand{\reto}{ 
		 {\Large RETOUR ET PROJECTION} \\
	\ding{117}	Dis ce que tu as appris ? \\
	\ding{117} Fait part de tes réussites puis de tes difficultés et la façon dont tu  les  as surmontées.\\
	\ding{117}	Dis ce en quoi ces résultats peuvent être utiles pour toi dans la vie courante\\
	\Ovalbox{\textbf{ Durée}:\dots }
}
\newcommand{\doc}{\begin{center}
		\section*{Documents utilisés}
	\end{center}
	\ding{117}	Guide et programme de la classe de troisième;\\
	\ding{117}	Réussir en mathématiques troisième;\\
	\ding{117} 	CIAM troisième.}	
\newcommand{\sph}{
 \begin{tikzpicture}[scale=0.7]
	\draw[color=red] (0,0) circle (1cm);
	\draw[dashed, color=red] (1,0) arc (0:180:1cm and 0.5cm);
	\draw[color=red] (-1,0) arc (180:360:1cm and 0.5cm);
	\draw[color=red] (0,-1) arc (270:90:0.5cm and 1cm);
	\draw[dashed, color=red!50] (0,1) arc (270:90:-0.5cm and -1cm);
\end{tikzpicture}}
\newcommand{\rect}{
\begin{tikzpicture}
	\draw[color=blue] (0,0)rectangle (0.7,0.7);
	\draw [color=blue] (0.7,0.7)--(1.5,1.5)--(1.5,0.8)--(0.7,0);
	\draw[color=blue] (1.5,1.5)--(0.8,1.5)--(0,0.7) ;
\end{tikzpicture}}
\newcommand{\cone}{
\begin{tikzpicture}[scale=0.7]
	\draw[dashed, color=orange] (0,0) arc (0:180:1cm and 0.5cm);
	\draw[color=orange] (-2,0) arc (0:180:-1cm and -0.5cm);
	\draw[color=orange](-2,0)--(-1,2)--(0,0);
\end{tikzpicture}}
\newcommand{\rapot}{
		%\draw (0,0) circle (5cm);
		%	\draw (0,0)--(5:5) node[right, font=\scriptsize] {5};
		\draw[line width=2pt] (0,0)--(0:5) node[right, font=\scriptsize]  {0};
		\draw (0,0)--(5:5) node[right, font=\scriptsize] {5} ;
		\draw (0,0)--(10:5) node[right, font=\scriptsize] {10};
		\draw (0,0)--(15:5) node[right, font=\scriptsize]{15} ;
		\draw (0,0)--(20:5) node[right, font=\scriptsize]{20} ;
		\draw (0,0)--(25:5) node[right, font=\scriptsize]{25} ;
		\draw (0,0)--(30:5) node[right, font=\scriptsize] {30};
		\draw (0,0)--(35:5) node[right, font=\scriptsize] {35};
		\draw (0,0)--(40:5) node[right, font=\scriptsize]{40} ;
		\draw (0,0)--(45.0:5) node[above right, font=\scriptsize]{45} ;
		\draw (0,0)--(50:5)node[above right, font=\scriptsize]{50} ;
		\draw (0,0)--(55.0:5)node[above right, font=\scriptsize]{55} ;
		\draw (0,0)--(60:5)node[above right, font=\scriptsize] {60};
		\draw (0,0)--(65.0:5)node[above right, font=\scriptsize] {65};
		\draw (0,0)--(70:5)node[above right, font=\scriptsize]{70} ;
		\draw (0,0)--(75.0:5) node[above, font=\scriptsize]{75};
		\draw (0,0)--(80:5)node[above , font=\scriptsize]{80};
		\draw (0,0)--(85:5)node[above, font=\scriptsize]{85};
		\draw[line width=2pt] (0,0)--(90:5)node[above, font=\scriptsize]{90};
		\draw (0,0)--(175:5) node[left, font=\scriptsize] {175} ;
		\draw (0,0)--(170:5) node[left, font=\scriptsize] {170};
		\draw (0,0)--(165:5) node[left, font=\scriptsize]{165} ;
		\draw (0,0)--(160:5) node[left, font=\scriptsize]{160} ;
		\draw (0,0)--(155:5) node[left, font=\scriptsize]{155} ;
		\draw (0,0)--(150:5) node[left, font=\scriptsize] {150};
		\draw (0,0)--(145:5) node[left, font=\scriptsize] {145};
		\draw (0,0)--(140:5) node[left, font=\scriptsize]{140} ;
		\draw (0,0)--(135.0:5) node[above left, font=\scriptsize]{135} ;
		\draw (0,0)--(130:5)node[above left, font=\scriptsize]{130} ;
		\draw (0,0)--(125.0:5)node[above left, font=\scriptsize]{125} ;
		\draw (0,0)--(120:5)node[above left, font=\scriptsize] {120};
		\draw (0,0)--(115.0:5)node[above left, font=\scriptsize] {115};
		\draw (0,0)--(110:5)node[above left, font=\scriptsize]{110} ;
		\draw (0,0)--(105.0:5) node[above, font=\scriptsize]{105};
		\draw (0,0)--(100:5)node[above , font=\scriptsize]{100};
		\draw (0,0)--(95:5)node[above, font=\scriptsize]{95};
		\draw [line width=2pt](0,0)--(180.0:5) node[left, font=\scriptsize] {180};
		\draw [fill=white](4,0) arc (0:180:4);
		\draw [](5,0) arc (0:180:5);
		\draw [line width=2pt] (0,0)--(90:4);
		\draw [line width=2pt] (0,0) rectangle (0.2,0.2);
}
\newcommand{\pavdroi}{ 
		\draw  (0,0)rectangle (4,2);
		\draw  (4.7,2.7)--(0.7,2.7);
		\draw[dashed,] (0,0)--(0.7,0.7);
		\draw  [color=black] (4,0)--(4.7,0.7) ;
		\draw  (0,2)--(0.7,2.7);
		\draw  (4,02)--(4.7,2.7) ;
		\draw [dashed] (0.7,0.7)--(0.7,2.7);
		\draw [dashed] (0.7,0.7)--(4.7,0.7);
		\draw  (0,0)--(0.7,0.7)--(4.7,0.7)--(4,0)--cycle;
		\draw  (0,2)--(0.7,2.7)--(4.7,2.7)--(4,2)--cycle;
		\draw  (4.7,0.7)-- (4.7,2.7) ;
		\draw (0,0)--(4,0);
}
\newcommand{\cube}{\begin{tikzpicture}[x={(-0.572cm,-0.416cm)},y={(1cm,0cm)},z={(0cm,1cm)}]
		\draw  (0,0,0)--(1,0,0)--(1,1,0)--(0,1,0)--cycle;
		\draw   (0,0,1)--(1,0,1)--(1,1,1)--(0,1,1)--cycle; 
		\draw (0,0,0)--(0,0,1);
		\draw (1,0,0)--(1,0,1);
		\draw (0,1,0)--(0,1,1); 
		\draw (1,1,0)--(1,1,1) ; 
		\draw  (1,0,0)--(1,1,0) ;
		\draw  (1,1,0)--(0,1,0);
	
\end{tikzpicture}}
%\SetWatermarkText{KAYCT}